% =====================================================================
% Car vs Truck - resultados por classe nas 4 configuracoes
% Fonte: paper/repro/car-truck-breakdown.json (== paper/pillar23/pillar3-single-classfix.json)
% Gerado em 2026-09-07. Valores EXATOS, nao arredondados aqui.
%
% Definicao de "correct" (identica nas 4 linhas, sem IoU contra o gabarito):
%   a caixa sobrevive ao estagio de curadoria E
%   (veredito humano in {kept, corrected-minor} com o rotulo preservado)
%   OU (veredito humano == class-changed e o rotulo final == o que o humano pos).
%   O vinculo proposta<->gabarito e a identidade ref_det_id, nao sobreposicao.
%
% Denominadores (conferidos):
%   car   : GT 427 = 133 never proposed + 294 with a proposal ; 304 proposals
%   truck : GT 255 =  93 never proposed + 162 with a proposal ; 149 proposals
%   coverage         = correct / GT da classe
%   reachable recov. = correct / GT da classe que TINHA proposta
%   not recovered    = GT da classe - correct
% =====================================================================

% ---------------------------------------------------------------------
% TABELA COMPACTA (recomendada para o corpo do paper)
% ---------------------------------------------------------------------
\begin{table}[t]
\centering
\caption{Per-class outcome of each curation stage on the 313-frame reference
set. \emph{Boxes} is how many detector proposals survive the stage;
\emph{correct} follows the verdict-based definition (no IoU matching against the
reference; proposals are linked to reviewer actions by identity). Coverage is
over all reference instances of the class; reachable recovery is over the
instances the detector actually proposed. Class-correction counts are
attempted/correct/incorrect.}
\label{tab:cartruck}
\small
\setlength{\tabcolsep}{4pt}
\begin{tabular}{llrrrrrrr}
\toprule
Class & Stage & Boxes & Corr. & Inc. & Not rec. & Prec. & Cov. & Reach. \\
\midrule
\multirow{4}{*}{\shortstack[l]{Car\\\scriptsize GT 427\\\scriptsize 294 proposed}}
 & YOLO11x (no VLM)          & 304 & 208 & 96 & 219 & 68.4 & 48.7 & 70.7 \\
 & VLM single (v1)           & 255 & 192 & 63 & 235 & 75.3 & 45.0 & 65.3 \\
 & VLM single + class (v3)   & 255 & 192 & 63 & 235 & 75.3 & 45.0 & 65.3 \\
 & VLM decomposed (v2)       & 283 & \textbf{217} & 66 & 210 & \textbf{76.7} & \textbf{50.8} & \textbf{73.8} \\
\midrule
\multirow{4}{*}{\shortstack[l]{Truck\\\scriptsize GT 255\\\scriptsize 162 proposed}}
 & YOLO11x (no VLM)          & 149 & 107 & 42 & 148 & 71.8 & 42.0 & 66.0 \\
 & VLM single (v1)           & 139 & 107 & 32 & 148 & 77.0 & 42.0 & 66.0 \\
 & VLM single + class (v3)   & 139 & 110 & 29 & 145 & 79.1 & 43.1 & 67.9 \\
 & VLM decomposed (v2)       & 142 & \textbf{118} & 24 & 137 & \textbf{83.1} & \textbf{46.3} & \textbf{72.8} \\
\bottomrule
\end{tabular}
\end{table}

% ---------------------------------------------------------------------
% TABELA COMPLETA (com correcao de classe; para o appendix)
% ---------------------------------------------------------------------
\begin{table}[t]
\centering
\caption{Per-class outcome with class-correction accounting. Denominators:
car, 427 reference instances of which 294 received a detector proposal and 133
were drawn by hand; truck, 255 of which 162 received a proposal and 93 were
drawn by hand. The detector emitted 304 car and 149 truck proposals. Rows are
indexed by the class the \emph{detector} proposed, so a class correction moves
the credit to the proposed-class row, not the reference-class row.}
\label{tab:cartruck-full}
\scriptsize
\setlength{\tabcolsep}{3.5pt}
\begin{tabular}{llrrrrrrrrrr}
\toprule
 & & \multicolumn{4}{c}{Counts} & \multicolumn{3}{c}{Rates (\%)} & \multicolumn{3}{c}{Class corrections} \\
\cmidrule(lr){3-6}\cmidrule(lr){7-9}\cmidrule(lr){10-12}
Class & Stage & Boxes & Corr. & Inc. & Not rec. & Prec. & Cov. & Reach. & Att. & OK & Bad \\
\midrule
\multirow{4}{*}{Car}
 & YOLO11x                 & 304 & 208 & 96 & 219 & 68.4 & 48.7 & 70.7 &  0 & 0 & 0 \\
 & VLM single (v1)         & 255 & 192 & 63 & 235 & 75.3 & 45.0 & 65.3 &  1 & 0 & 1 \\
 & VLM single + class (v3) & 255 & 192 & 63 & 235 & 75.3 & 45.0 & 65.3 &  1 & 1 & 0 \\
 & VLM decomposed (v2)     & 283 & 217 & 66 & 210 & 76.7 & 50.8 & 73.8 & 11 & 9 & 0 \\
\midrule
\multirow{4}{*}{Truck}
 & YOLO11x                 & 149 & 107 & 42 & 148 & 71.8 & 42.0 & 66.0 &  0 & 0 & 0 \\
 & VLM single (v1)         & 139 & 107 & 32 & 148 & 77.0 & 42.0 & 66.0 &  5 & 0 & 4 \\
 & VLM single + class (v3) & 139 & 110 & 29 & 145 & 79.1 & 43.1 & 67.9 &  5 & 4 & 1 \\
 & VLM decomposed (v2)     & 142 & 118 & 24 & 137 & 83.1 & 46.3 & 72.8 & 15 &12 & 1 \\
\bottomrule
\end{tabular}
\end{table}

% ---------------------------------------------------------------------
% NOTA DE RODAPE (usar uma das duas)
% ---------------------------------------------------------------------
% \footnote{Rows are indexed by the class the detector proposed. Because a class
% correction moves the credit to the proposed-class row rather than to the
% reference-class row, per-class coverage is only interpretable where the two
% partitions nearly coincide, which holds for car and truck. The remaining
% classes (bus 23, motorcycle 13, person 10, bicycle 3 reference instances) are
% too small to support a per-class claim and are omitted.}

% \footnote{Attempted class corrections need not equal correct plus incorrect:
% a correction issued on a proposal whose reviewer verdict was a geometry
% redraw is counted as attempted but is neither credited nor charged.}

% ---------------------------------------------------------------------
% RECONCILIACAO COM A TABELA PRINCIPAL (conferido em 2026-09-07)
% Somar as linhas por classe e comparar com a tabela principal NAO bate,
% e a diferenca e conhecida e inteira:
%
%   caixas retidas    base: 499 por classe  vs  501 na principal  (-2)
%                       v2: 467 por classe  vs  468 na principal  (-1)
%   correcoes classe    v2:  53 tentadas /  42 certas por classe
%                            54 tentadas /  43 certas na principal (-1 / -1)
%
% Causa unica: o recorte por classe so emite linha para classe que EXISTE no
% gabarito, e duas das 501 propostas tem classe que o gabarito nao tem --
% uma `fork` (o humano rejeitou; sai do v2 tambem) e uma `fire hydrant`
% (o humano trocou para motorcycle e o v2 acertou a troca -- e essa a 54a
% correcao tentada e a 43a certa da tabela principal).
% Nenhuma das duas e car nem truck: as linhas de car e de truck fecham
% EXATAMENTE com a tabela principal, e os numeros acima nao as tocam.
% Fonte: car-truck-breakdown.json + phase1.json do run do v2.
